\documentclass[12pt]{article}
\usepackage[T1]{fontenc}
\usepackage[utf8]{inputenc}
\usepackage{lmodern}
\usepackage{textcomp}
\usepackage{enumitem}
\usepackage{geometry}
\usepackage{graphicx}
\usepackage{amsmath, amssymb}
\geometry{margin=1in}
\begin{document}
\begin{center}
History - Graduate Level Assessment 14 \
Total Marks: 40 \
Answer all questions.
\end{center}

\section*{Multiple Choice (10 marks)}
\begin{enumerate}[label=\arabic*.]
\item The possibility that Homo erectus possessed an aesthetic sense and may have created "art" is suggested by:
\begin{enumerate}[label=(\alph*)]
    \item the discovery of figurines made from clay at the Trinil site.
    \item recovery of a shell at the Trinil site with an intentionally engraved series of lines.
    \item the use of complex tools requiring a high level of cognitive function.
    \item cave paintings and elaborate burials with numerous grave goods at the Trinil site.
    \item all of the above.
\end{enumerate}
\item The "Lion Man" from Hohlenstein-Stadel cave is an example of:
\begin{enumerate}[label=(\alph*)]
    \item mobiliary art.
    \item long-distance trade of exotic raw materials.
    \item parietal art.
    \item a Venus figurine.
\end{enumerate}
\item When the Spanish explorer Hernando de Soto crossed a portion of the North American continent, his expedition encountered direct descendants of which people?
\begin{enumerate}[label=(\alph*)]
    \item Mayan
    \item Kwakiutl
    \item Mississippian
    \item Inca
    \item Clovis
\end{enumerate}
\item The skulls of the earliest migrants to the New World do not match those of modern Native Americans. This suggests that:
\begin{enumerate}[label=(\alph*)]
    \item the early migrants were not human.
    \item the early migrants did not contribute to the genetic pool of modern Native Americans.
    \item the early migrants were a separate species.
    \item craniometrics is an unreliable method of tracking populations.
    \item skull shapes can change dramatically over generations.
\end{enumerate}
\item Taphonomy is:
\begin{enumerate}[label=(\alph*)]
    \item the study of plant and animal life from a particular period in history.
    \item a study of the spatial relationships among artifacts, ecofacts, and features.
    \item the study of human behavior and culture.
    \item a technique to determine the age of artifacts based on their position in the soil.
    \item a dating technique based on the decay of a radioactive isotope of bone.
\end{enumerate}
\end{enumerate}

\section*{True / False (10 marks)}
\textit{Answer True or False}
\begin{enumerate}[label=\arabic*.]
\setcounter{enumi}{5}
\item True or False: Kievan Rus' begins with the rule (882-912) of Prince Oleg, who extended his control from Novgorod south along the Dnieper river valley in order to protect trade from Khazar incursions from the east and moved his capital to the more strategic Kiev. Sviatoslav I (died 972) achieved the first major expansion of Kievan Rus' territorial control, fighting a war of conquest against the Khazar Empire.
\item True or False: Most other member states have expressed similar reservations. This suggests the EU's legitimacy rests on the ultimate authority of member states, its factual commitment to human rights, and the democratic will of the people.
\item True or False: In 1928, Nasser went to Alexandria to live with his maternal grandfather and attend the city's Attarin elementary school.
\item True or False: Many of them were nephews, resulting in Carl, Sr. being known around the studios as "Uncle Carl." Ogden Nash famously quipped in rhyme, "Uncle Carl Laemmle/Has a very large faemmle." Among these relatives was future Academy Award winning director/producer William Wyler.
\item True or False: The conflict lasted one week before both sides agreed to sign a peace treaty that was brokered by several Arab states. That year, Gaddafi was invited to Moscow by the Soviet government in recognition of their increasing a different answer relationship.
\end{enumerate}

\section*{Long Form Response (20 marks)}
\textit{Answer each question with a clear and complete explanation.}
\begin{enumerate}[label=\arabic*.]
\setcounter{enumi}{10}
\item Read the following passage and answer the question: The study of genocide has mainly been focused towards the legal aspect of the term. By formally recognizing the act of genocide as a crime, involves the undergoing prosecution that begins with not only seeing genocide as outrageous past any moral standpoint but also may be a legal liability within international relations. When genocide is looked at in a general aspect it is viewed as the deliberate killing of a certain group. Yet is commonly seen to escape the process of trial and prosecution due to the fact that genocide is more often than not committed by the officials in power of a state or area. In 1648 before the term genocide had been coined, the Peace of Westphalia was established to protect ethnic, national, racial and in some instances religious groups. During the 19 th century hu... Question: In prosecuting genocide, what must the act be formally acknowledged as?
\item Read the following passage and answer the question: Georgian succeeded the English Baroque of Sir Christopher Wren, Sir John Vanbrugh, Thomas Archer, William Talman, and Nicholas Hawksmoor; this in fact continued into at least the 1720 s, overlapping with a more restrained Georgian style. The architect James Gibbs was a transitional figure, his earlier buildings are Baroque, reflecting the time he spent in Rome in the early 18 th century, but he adjusted his style after 1720. Major architects to promote the change in direction from baroque were Colen Campbell, author of the influential book Vitruvius Britannicus (1715-1725); Richard Boyle, 3 rd Earl of Burlington and his protégé William Kent; Isaac Ware; Henry Flitcroft and the Venetian Giacomo Leoni, who spent most of his career in England. Other prominent architects of the early Georgian ... Question: What Venetian spent most of his career in England?
\end{enumerate}

\vfill
\noindent\textit{End of Paper}
\end{document}